% Study sheet: Volume IV, Chapter 2, The Physical Body.
% Numbers: \val{...} from body_studysheet_numbers.tex, written by body_studysheet_numbers.py.
\documentclass{studysheet}
\input{body_studysheet_numbers}
\sheetvolume{4}
\sheetchapter{2}{The Physical Body}
\sheetsources{Numbers from the chapter's mlsysim scenarios for the warehouse mobile manipulator (illustrative profiles); questions marked assumed use new inputs.}

\begin{document}
% ===================================================================== FRONT: LEARN
\frontside
\begin{bigidea}
A correct decision can still ask more than the body can deliver. Before a command becomes current
in a winding, the permission path checks \hl{five physical budgets at once}. Work across the
\hl{causal boundary} cannot be recalled, so an overdrawn budget means refusal.
\end{bigidea}

\begin{leftcol}
\sheetsec{Key ideas}
\begin{ideas}
\idea{Five budgets, one admission.}{Slack in four never pays the fifth's deficit.}
\idea{Age runs from capture to onset.}{Inference is one stage of eight. The end-to-end tail
  decides, and stage tails do not add.}
\idea{Credible deceleration is the lowest ceiling.}{Traction, tip-over and drive torque, owned by the
  floor, the arm's pose and the drive settings.}
\idea{Gearing trades torque for inertia.}{Ratio $\GearReductionRatio$ multiplies torque by
  $\GearReductionRatio$, rotor inertia by $\GearReductionRatio^2$; for one motor, heat times inertia is fixed.}
\idea{Torque is current; current is heat.}{Loss grows with torque squared; thermal history spends
  the margin that peak torque needs.}
\idea{A limit is a record, not a number.}{Value, conditions, uncertainty, margin, detector,
  provenance. Only measured or measurement-derived values support a margin.}
\end{ideas}

\vspace{9pt}
\sheetsec{The admission check}
\begin{sheettable}{@{}l l L@{}}
\thead{Budget} & \thead{Must hold} & \thead{What spends it}\\ \midrule
Freshness & $\ObservationAge \le \TelemetryFreshnessBound$ & every stage to torque onset\\
Stopping & $\StoppingDistance \le \ClearDistance$ & speed squared, delay, floor\\
Torque & $|\CommandedJointTorque| \le \OutputTorqueCeiling$ & current limit, back-EMF\\
Heat & $\EstimatedWindingTemperature \le \WindingTemperatureLimit$ & current squared, history\\
Supply & $\MinimumRailVoltage \ge \ConverterLockoutVoltage$ & coincident load, pack age\\
\end{sheettable}

\vspace{9pt}
\fcard{Shared rail}%
  {\SourceVoltage - \FiniteChange\ConductorCurrent\,\PowerRailResistance \ge \ConverterLockoutVoltage,
   \qquad \FinalBusVoltage = \sqrt{\InitialBusVoltage^2 + \frac{2\RegeneratedBrakingEnergy}{\DcBusBulkCapacitance}}}%
  {Sag under a coincident step $\FiniteChange\ConductorCurrent$; bus after a stop returns $\RegeneratedBrakingEnergy$
   into capacitance alone. Past a switch rating, fit a brake chopper.}%
  {Sustained only; transients dip lower.}{}
\end{leftcol}%
\begin{rightcol}
\fcard{Stopping distance and speed ceiling}%
  {\begin{gathered}\StoppingDistance = \SpeedAtDetection\PrebrakeDelay + \frac{\SpeedAtDetection^2}{2\CredibleBrakingDeceleration}
     + \LocalizationBound + \ProtectiveClearanceMargin \le \ClearDistance\\[2pt]
   \SpeedBound = -\CredibleBrakingDeceleration\PrebrakeDelay + \sqrt{(\CredibleBrakingDeceleration\PrebrakeDelay)^2
     + 2\CredibleBrakingDeceleration\AvailableBrakingDistance}\end{gathered}}%
  {Blind travel over pre-brake delay $\PrebrakeDelay$, braking, two fixed allowances;
   $\AvailableBrakingDistance = \ClearDistance - \LocalizationBound - \ProtectiveClearanceMargin$.}%
  {Delay enters linearly, speed quadratically.}{}
\fcard{Credible deceleration}%
  {\CredibleBrakingDeceleration \le \min\!\left(\CoulombFrictionCoefficient\Gravity\DrivenWheelLoadFraction,\;
     \frac{\Gravity\LeadingSupportDistance}{\GravityCenterHeight},\;
     \frac{2\MotorTorqueConstant\InverterCurrentLimit}{\WheelRadius\,\EffectiveImpactMass}\right)}%
  {Traction (driven-wheel load share $\DrivenWheelLoadFraction$), tip-over, two hub drives at current limit
   $\InverterCurrentLimit$ ($\MotorTorqueConstant$: torque constant; $\EffectiveImpactMass$: effective mass). A drive can clamp
   yet acknowledge: watch measured phase current and its saturation flag.}%
  {}%
  {\begin{tikzpicture}[sheet]
     \pgfmathsetmacro\sc{32}
     \foreach \lab/\v/\y/\sty in {{dry, rated share}/\val{adry}/45/nbox, {drives, peak current}/\val{adrive}/30/nbox, {inspected, all weight}/\val{ainsp}/15/hotbox}{
       \node[lab, anchor=east] at (78,\y+5) {\lab};
       \draw[\sty] (82,\y) rectangle ({82+\v*\sc},\y+10);
       \node[note, anchor=west] at ({85+\v*\sc},\y+5) {\v};}
     \node[lab, anchor=east] at (78,5) {tip-over};
     \node[note, anchor=west] at (85,5) {not measured};
     \draw[dim, dashed] ({82+\val{areq}*\sc},-1) -- ({82+\val{areq}*\sc},58);
     \node[small, anchor=south] at ({82+\val{areq}*\sc},58) {needed \val{areq}};
     \node[note, anchor=north west] at (82,-3) {ceilings, m/s$^2$};
   \end{tikzpicture}}
\fcard{Freshness deadline}%
  {\TelemetryFreshnessBound \le \frac{\PositionErrorCeiling}{\ConveyorSpeed}}%
  {$\PositionErrorCeiling$: motion allowance (tolerance minus initial estimate error); $\ConveyorSpeed$: target speed.}%
  {Set by the task: a faster belt shrinks it with no change in the pipeline.}{}
\fcard{Burst horizon}%
  {\BurstSurvivalTime = \ThermalTimeConstant \ln\frac{\EquilibriumWindingTemperature - \WindingTemperature(0)}{\EquilibriumWindingTemperature - \WindingTemperatureLimit},
   \quad \EquilibriumWindingTemperature = \AmbientTemperature + \PeakCopperLossPower\ThermalResistance}%
  {Time for loss $\PeakCopperLossPower$ ($\WindingLossPower = \TorqueProducingCurrent^2\EffectiveLossResistance$) to carry the
   winding from $\WindingTemperature(0)$ to its insulation limit.}%
  {A case thermistor lags the winding; admit on the observer's estimate.}{}
\end{rightcol}

% ===================================================================== BACK: PRACTICE
\backside
\begin{answers}
\ans{Finish}{Total \val{wblind} + \val{ybrake} + \val{wover} = \res{\val{ytotal}\,mm}, \val{yover}\,mm past the
  \val{wclearmm}\,mm: refuse. With \val{ainsp} in the ceiling, $\SpeedBound \approx$ \res{\val{yvmax}\,m/s}, itself
  an upper bound (all weight on the driven wheels).}
\ans{1}{\val{emax}\,mm $\div$ \val{qvconv}\,m/s = \res{\val{qfresh}\,ms}. Median fits by \val{qmedspare}\,ms; p99 misses by
  \val{qpnnover}\,ms: no. Slow the belt or cut the tail; the backbone is not the problem (freshness card).}
\ans{2}{\val{tauth}\,s $\times \ln(\val{qnum}/\val{qden})$ = \res{\val{qhot}\,s} to the insulation limit (\val{qcold}\,s from ambient). The observer:
  the case lags the hotspot and overstates the margin (burst card).}
\columnbreak
\ans{3}{(a) \val{di}\,A $\times$ \val{qrbus}\,m$\Omega$ = \val{qsag}\,V; \val{qvbus}\,V is only \res{\val{qmargin}\,V} above lockout,
  before any transient dip. (b) About \res{\val{qvfin}\,V}, above \val{vswitch}\,V: keep the chopper (rail card).}
\ans{4}{Conditions (holds only at \val{ipeak}\,A). Twice \val{kt} $\times$ \val{irated} $\div$ \val{rwheel}\,mm $\approx$ \val{qforce}\,N;
  $\div$ \val{meff}\,kg = \res{\val{qadrive}\,m/s²}, below the \val{areq} needed. Detector: measured phase current
  and saturation flag, not the acknowledgement (record).}
\end{answers}

\begin{leftcol}
\sheetsec{Worked example}
\begin{practicebox}\boxtitle{Can the loaded base stop before the rack end?}
The base runs at its \val{wv}\,m/s drive limit on a dry floor ($\CredibleBrakingDeceleration$ = \val{areq}\,m/s$^2$).
A person stands \val{wclear}\,m ahead. Full pre-brake delay \val{wtau}\,ms; localization bound \val{wloc}\,mm
and protective clearance \val{wmar}\,mm.\par\vspace{5pt}
{\centering
\begin{tikzpicture}[sheet]
  \draw[hotbox] (0,0) rectangle (\val{pblind},12);
  \draw[accbox] (\val{pblind},0) rectangle ({\val{pblind}+\val{pbrake}},12);
  \draw[nbox] ({\val{pblind}+\val{pbrake}},0) rectangle ({\val{pblind}+\val{pbrake}+\val{pover}},12);
  \node[small, anchor=south] at ({\val{pblind}/2},13) {blind};
  \node[small, anchor=south] at ({\val{pblind}+\val{pbrake}/2},13) {braking};
  \node[small, anchor=south] at ({\val{pblind}+\val{pbrake}+\val{pover}/2},13) {overhead};
  \draw[dim, dashed] (\val{pclear},-4) -- (\val{pclear},24);
  \node[small, anchor=south] at (\val{pclear},24) {rack end};
  \node[note, anchor=north] at (110,-4) {to scale; spare \val{wspare}\,mm};
\end{tikzpicture}\par}\vspace{4pt}
\begin{steps}
\item Blind travel: \val{wv}\,m/s $\times$ \val{wtau}\,ms = \res{\val{wblind}\,mm}.
\item Braking: $\val{wv}^2 / (2\times\val{areq})$ = \res{\val{wbrake}\,mm}.
\item Plus \val{wover}\,mm overhead: \res{\val{wtotal}\,mm}, inside \val{wclearmm}\,mm, if the deceleration holds.
\item Ceiling with $\AvailableBrakingDistance$ = \val{wavail}\,m: $\SpeedBound \approx$ \res{\val{wvmax}\,m/s}.
\end{steps}
\takeaway{The model fits; admission waits on evidence for its premise: loaded wheel share and tip-over are unmeasured.}
\end{practicebox}

\vspace{10pt}
\sheetsec{Your turn: finish it}
\begin{practicebox}\boxtitle{Same stop, inspected floor}
The aisle is the inspected floor ($\CoulombFrictionCoefficient$ = \val{muinsp}). Even with all weight on the
driven wheels, traction caps $\CredibleBrakingDeceleration$ at \val{ainsp}\,m/s$^2$.\par\vspace{4pt}
\textbf{Given:} blind travel is still \val{wblind}\,mm; braking is $\val{wv}^2/(2\times\val{ainsp})$ = \val{ybrake}\,mm.\par
\textbf{You:} find the total, decide admit or refuse, and find the new speed ceiling.\par\vspace{30pt}
\end{practicebox}
\end{leftcol}%
\begin{rightcol}
\sheetsec{Common mistakes}
\mistake{``The backbone takes \val{backbone}\,ms against a \val{fresh}\,ms deadline, so the grasp is fresh.''}%
  {Age runs from exposure to torque onset: \val{agemed}\,ms at the median and \val{agetail}\,ms at p99.9, when
   the mug moves \val{traveltail}\,mm against a \val{emax}\,mm allowance.}{}
\mistake{``A higher gear ratio is free: the motor runs cooler.''}%
  {From \val{gone}:1 to \val{gtwo}:1 the arm joint's holding loss falls from \val{lossone} to \val{losstwo}\,W, but
   its reflected inertia rises from \val{jone} to \val{jtwo}\,kg$\cdot$m$^2$, which an impact must stop. Only a better
   motor improves both.}{}

\vspace{2pt}
\sheetsec{Check yourself}
\begin{qs}
\item \assumed{} the belt runs at \val{qvconv}\,m/s; the motion allowance stays \val{emax}\,mm and the ages
  (median \val{agemed}\,ms, p99 \val{agepnn}\,ms) are unchanged. Deadline? Does it meet a p99 requirement?
\item \assumed{} after a hold the observer estimates \val{qtzero}\,\textdegree C; the case reads \val{qcase}\,\textdegree C.
  A \val{ppk}\,W burst, $\ThermalResistance$ = \val{rth}\,K/W, $\ThermalTimeConstant$ = \val{tauth}\,s, ambient
  \val{tamb}\,\textdegree C, limit \val{tlim}\,\textdegree C: time to the limit? Which reading counts?
\item \assumed{} an aged, cold pack: $\PowerRailResistance$ = \val{qrbus}\,m$\Omega$. (a) Margin for a \val{di}\,A step
  from \val{vlow}\,V, lockout \val{vuvlo}\,V? (b) To drop the chopper, $\DcBusBulkCapacitance$ doubles to
  \val{qcbus}\,$\mu$F; a stop returns \val{eregen}\,J into a \val{vfull}\,V bus, switches rated \val{vswitch}\,V. Safe?
\item To protect the windings, the hub drives' limit drops from \val{ipeak} to \val{irated}\,A
  ($\MotorTorqueConstant$ = \val{kt}\,N$\cdot$m/A, $\WheelRadius$ = \val{rwheel}\,mm, $\EffectiveImpactMass$ = \val{meff}\,kg). Which
  record element breaks? New drive ceiling? What detects it?
\end{qs}
\end{rightcol}
\end{document}
